\documentclass[12pt,oneside]{book}
\usepackage[margin=1in]{geometry}
\usepackage{iftex}
\ifPDFTeX
  \usepackage[T1]{fontenc}
  \usepackage[utf8]{inputenc}
  \usepackage{lmodern}
\else
  \usepackage{fontspec}
  \setmainfont{Latin Modern Roman}
\fi
\usepackage{amsmath,amssymb,amsthm,mathtools}
\usepackage{longtable,booktabs}
\usepackage{enumitem}
\usepackage{setspace}
\usepackage{fancyhdr}
\usepackage{titlesec}
\usepackage[dvipsnames]{xcolor}
\usepackage{xurl}
\usepackage{textcomp}
\usepackage{hyperref}
\usepackage[nameinlink,noabbrev]{cleveref}
\usepackage[most]{tcolorbox}
\usepackage{microtype}
\usepackage{framed}
\usepackage{fancyvrb}

\hypersetup{
  colorlinks=true,
  linkcolor=blue!50!black,
  citecolor=blue!50!black,
  urlcolor=blue!60!black,
  pdftitle={The Four Canaries: UT Austin Prelim Problems in Lean~4},
  pdfauthor={Joseph Corneli and Claude}
}

\setcounter{tocdepth}{2}
\setcounter{secnumdepth}{3}
\setlength{\parskip}{0.2em}
\setlength{\parindent}{0pt}
\setlist[itemize]{itemsep=0.3em,topsep=0.3em}
\setlist[enumerate]{itemsep=0.3em,topsep=0.3em}
\setstretch{1.08}

\theoremstyle{plain}
\newtheorem{theorem}{Theorem}[chapter]
\newtheorem{lemma}[theorem]{Lemma}
\newtheorem{corollary}[theorem]{Corollary}
\theoremstyle{definition}
\newtheorem{problem}{Problem}[chapter]
\theoremstyle{remark}
\newtheorem*{remark}{Remark}

\providecommand{\tightlist}{}
\definecolor{shadecolor}{RGB}{248,248,248}
\DefineVerbatimEnvironment{Highlighting}{Verbatim}{commandchars=\\\{\},fontsize=\small}
\newenvironment{Shaded}{\begin{snugshade}}{\end{snugshade}}

% Lean code box
\newtcolorbox{leanbox}[1][]{%
  colback=blue!2, colframe=blue!30!black,
  fonttitle=\bfseries\sffamily\small,
  title={Lean~4 / Mathlib}, #1,
  left=6pt, right=6pt, top=4pt, bottom=4pt,
  boxrule=0.5pt, arc=2pt,
  before skip=6pt, after skip=6pt
}

% Informal proof box
\newtcolorbox{proofbox}[1][]{%
  colback=orange!3, colframe=orange!40!black,
  fonttitle=\bfseries\sffamily\small,
  title={Informal Proof}, #1,
  left=6pt, right=6pt, top=4pt, bottom=4pt,
  boxrule=0.5pt, arc=2pt,
  before skip=6pt, after skip=6pt
}

% Note box
\newtcolorbox{notebox}[1][]{%
  colback=green!3, colframe=green!40!black,
  fonttitle=\bfseries\sffamily\small,
  title={Process Note}, #1,
  left=6pt, right=6pt, top=4pt, bottom=4pt,
  boxrule=0.5pt, arc=2pt,
  before skip=6pt, after skip=6pt
}

\newcommand{\RP}{\mathbb{R}\mathrm{P}}

\titleformat{\chapter}[display]
  {\normalfont\Large\bfseries}
  {\filright\large\chaptername\ \thechapter}
  {8pt}{\LARGE}
\titlespacing*{\chapter}{0pt}{-6pt}{18pt}

\begin{document}

\begin{titlepage}
\vspace*{3cm}
\begin{center}
{\Huge\bfseries The Four Canaries}\\[0.8cm]
{\Large UT Austin Preliminary Exam Problems\\Formalised in Lean~4 with Mathlib}\\[2cm]
{\large Joseph Corneli \quad and \quad Claude (Opus 4.6)}\\[0.5cm]
{\normalsize Hyperreal Enterprises}\\[2cm]
{\normalsize 29 March 2026}\\[3cm]
{\small Four graduate-level mathematics problems---topology, analysis,\\
algebra, and functional analysis---proved informally\\
and verified formally, in a single session.}
\end{center}
\end{titlepage}

\tableofcontents

% ====================================================================
\chapter*{Abstract}
\addcontentsline{toc}{chapter}{Abstract}

Four problems from the University of Texas at Austin mathematics
preliminary examinations (1990--2003) were selected as ``canary''
problems for the M-apm-solutions mission: one each from topology,
analysis, algebra, and functional analysis. Each was proved
informally, then formalised in Lean~4 against Mathlib. All four
compile without \texttt{sorry}.

The algebra canary (C3) is a one-screen proof. The topology canary
(C1) reduces a covering-space argument to a verified algebraic
lemma about torsion-free groups. The functional analysis canary (C4)
carries the full Closed Graph Theorem argument through 94~lines.
The analysis canary (C2) is the hardest: 356~lines assembling
Fatou's lemma, H\"older's inequality, uniform integrability,
Vitali convergence, and an $\varepsilon/3$ tail argument to prove
weak $L^2$ convergence.

The experience suggests that formalisation is not a bottleneck for
problems where Mathlib coverage is strong (algebra, functional
analysis), and that the real cost is API composition---wiring
existing lemmas through Lean's type system---rather than missing
mathematics.

% ====================================================================
\part{The Lessons}
% ====================================================================
% Part I: The Lessons
% Each canary as a tutorial, not a proof demonstration.
% Target reader: strong undergrad who needs tools explained.
% ====================================================================

% --------------------------------------------------------------------
\chapter{C1: Topology --- Why Non-Surjective Maps Kill Themselves}\label{ch:lesson-c1}

\section{The setup}

$\RP^2$ is the \textbf{real projective plane}. You build it by taking a
sphere and gluing opposite points together: every point~$p$ on~$S^2$ gets
identified with the point directly on the other side,~$-p$. The result
is a surface that's locally flat (like a plane) but globally has a twist
in it---a bit like a M\"obius strip, except it's a closed surface with no
boundary.

The fundamental group $\pi_1(\RP^2)$ is $\mathbb{Z}/2\mathbb{Z}$---the
group with just two elements, $\{0, 1\}$, where $1 + 1 = 0$. Intuitively:
there's one nontrivial loop in~$\RP^2$ (you walk from a point to its
antipode on the sphere, and the identification glues the endpoints
together), and if you walk that loop \emph{twice}, you've gone all the
way around the sphere and returned---so the loop squares to the identity.

The \textbf{universal cover} of $\RP^2$ is~$S^2$ itself. The covering map
$p : S^2 \to \RP^2$ sends each point to its equivalence class. This is a
$2$-to-$1$ map. The key property of covering spaces: a continuous map
$f : \RP^2 \to \RP^2$ that induces the \emph{trivial} map on~$\pi_1$ can
be ``lifted'' to a map $\tilde{f} : \RP^2 \to S^2$ satisfying
$p \circ \tilde{f} = f$.

\section{What the problem is really asking}

If $f$ misses a point, does that force $f$ to be ``topologically
trivial''---homotopic to a constant? The answer is yes, and the reason is
\emph{algebraic}: the factorisation through $\RP^2 \setminus \{q\}$
forces $f_*$ to be the zero map, and once $f_*$ is zero, the lifting
property finishes the job.

\section{The tool you need: what does $\RP^2$ minus a point look like?}

This is the hint. Take the sphere~$S^2$ and remove a point~$q$ and its
antipode~$-q$. You get $S^2 \setminus \{q, -q\}$, which is homeomorphic
to the cylinder $S^1 \times \mathbb{R}$ (think of it as the sphere with
the north and south poles removed---that's a cylinder). When you glue
antipodes, the cylinder becomes an \textbf{open M\"obius band}.

An open M\"obius band deformation-retracts onto its core circle~$S^1$.
So $\pi_1(\RP^2 \setminus \{q\}) \cong \pi_1(S^1) \cong \mathbb{Z}$.

\section{The key insight}

Since $f$ misses~$q$, it factors: $f = i \circ g$ where
$g : \RP^2 \to \RP^2 \setminus \{q\}$ and $i$ is the inclusion. On
fundamental groups:
\[
  f_* = i_* \circ g_* :
  \underbrace{\mathbb{Z}/2\mathbb{Z}}_{\pi_1(\RP^2)}
  \xrightarrow{g_*}
  \underbrace{\mathbb{Z}}_{\pi_1(\RP^2 \setminus \{q\})}
  \xrightarrow{i_*}
  \underbrace{\mathbb{Z}/2\mathbb{Z}}_{\pi_1(\RP^2)}.
\]

Now: \textbf{what is $g_*$?} It's a group homomorphism from
$\mathbb{Z}/2\mathbb{Z}$ to~$\mathbb{Z}$. There's only one: the zero
map. Why? If $\varphi(1) = k$, then $0 = \varphi(0) = \varphi(1 + 1) =
\varphi(1) + \varphi(1) = 2k$. In~$\mathbb{Z}$, $2k = 0$ forces $k = 0$.
(This is what ``$\mathbb{Z}$ is torsion-free'' means: no nonzero element
has finite order.)

So $g_* = 0$, hence $f_* = 0$. For part~(b): $f_*$ trivial means $f$
lifts to $\tilde{f} : \RP^2 \to S^2$. The lift also misses a point
(because $f$ does), so it factors through $S^2 \setminus \{p\} \cong
\mathbb{R}^2$. Since $\mathbb{R}^2$ is contractible (shrink everything
to the origin), $\tilde{f}$ is nullhomotopic, hence so is~$f$.

\section{Where you'd get stuck}

\begin{itemize}
\item \textbf{Trying to compute $\pi_1$ from scratch.} Don't. The exam
  gives you the hint precisely so you don't have to run van~Kampen or
  build a CW-complex. Just identify the space and read off $\pi_1$.
\item \textbf{Thinking $\RP^2 \setminus \{q\}$ is contractible.} It
  retracts to $S^1$, not to a point. The M\"obius band has a
  nontrivial core loop. Getting this wrong teaches you the difference
  between ``simple'' and ``contractible.''
\item \textbf{Trying to build the homotopy by hand.} The lifting
  property and contractibility of $\mathbb{R}^2$ do this for you. The
  proof is existential, not constructive.
\end{itemize}

% --------------------------------------------------------------------
\chapter{C2: Analysis --- Why Bounded Sequences Converge Weakly}\label{ch:lesson-c2}

\section{The setup}

$L^2(\mathbb{R})$ is the space of square-integrable functions on the real
line: functions~$f$ for which $\int_{-\infty}^{\infty} |f(x)|^2\,dx <
\infty$. The norm is $\|f\|_2 = (\int |f|^2)^{1/2}$. This is a Hilbert
space---a complete inner product space---with inner product $\langle f, g
\rangle = \int f(x)\,\overline{g(x)}\,dx$.

The closed unit ball $B = \{f : \|f\|_2 \le 1\}$ is the set of functions
with bounded energy.

``$f_n \to f$ almost everywhere'' means: for almost every $x \in
\mathbb{R}$, the sequence of numbers $f_n(x)$ converges to $f(x)$.
``Almost every'' means the set of exceptions has measure zero.

\section{The two tools you need}

\textbf{Fatou's lemma.} If $g_n \ge 0$ are measurable functions and
$g_n \to g$ pointwise a.e., then
\[
  \int g \le \liminf_{n \to \infty} \int g_n.
\]
In words: the integral of the limit is at most the limit of the
integrals. The inequality can be strict (mass can ``escape to infinity''),
but you never gain mass in the limit.

Applied to $g_n = |f_n|^2$: since $f_n \to f$ a.e., we have $|f_n|^2
\to |f|^2$ a.e., so
\[
  \|f\|_2^2 = \int |f|^2 \le \liminf_n \int |f_n|^2
  = \liminf_n \|f_n\|_2^2 \le 1.
\]
That's part~(i). One tool, one inequality, done.

\textbf{H\"older's inequality} (the case $p = q = 2$ is Cauchy--Schwarz).
If $f \in L^2$ and $g \in L^2$, then $fg \in L^1$ and
\[
  \left|\int fg\right| \le \|f\|_2 \cdot \|g\|_2.
\]
This will be the workhorse for part~(ii).

\section{What part (ii) is really asking}

The statement $\int f_n\,g \to \int f\,g$ for all $g \in L^2$ is called
\textbf{weak convergence} of $f_n$ to~$f$. It's weaker than
$\|f_n - f\|_2 \to 0$ (which is ``strong'' or ``norm'' convergence and is
NOT what's happening here---bounded sequences in $L^2$ can converge a.e.\
without converging in norm).

The exam is asking you to prove: \emph{almost-everywhere convergence plus
a uniform $L^2$ bound implies weak convergence.}

\section{The key insight: split and conquer}

The hint says ``start with $g$ continuous and compactly supported.'' But
even that special case has a trap: there's no \emph{pointwise} dominator
for $f_n \cdot g$ that works for all~$n$ (because $\|f_n\|_2 \le 1$
does NOT give you $|f_n(x)| \le M$ for any fixed~$M$). So the Dominated
Convergence Theorem doesn't apply directly.

The way through: \textbf{split the integral into a core and a tail}.

Fix $\varepsilon > 0$. Choose a set $K$ (think $[-R, R]$) large enough
that $g$ has very little $L^2$ mass outside $K$:
$\|g \cdot \mathbf{1}_{K^c}\|_2 < \varepsilon$. This is possible
because $g \in L^2$ and $L^2$ tails vanish.

\textbf{The tail piece} is easy by Cauchy--Schwarz:
\[
  \left|\int_{K^c} (f_n - f)\,g\right|
  \le \|f_n - f\|_2 \cdot \|g \cdot \mathbf{1}_{K^c}\|_2
  \le 2\varepsilon.
\]
(We used $\|f_n - f\|_2 \le \|f_n\|_2 + \|f\|_2 \le 1 + 1 = 2$.)

\textbf{The core piece} needs more work. On $K$ (a bounded interval, so
finite measure), we can use:

\begin{enumerate}
\item \textbf{Uniform integrability} of the product $(f_n - f) \cdot g$
  on~$K$. The estimate: for any small set $A \subseteq K$,
  \[
    \|\mathbf{1}_A \cdot (f_n - f) \cdot g\|_1
    \le \|f_n - f\|_2 \cdot \|\mathbf{1}_A \cdot g\|_2
    \le 2 \cdot \|\mathbf{1}_A \cdot g\|_2.
  \]
  The right side goes to zero as $\mu(A) \to 0$ (because $g \in L^2$
  and the $L^2$ norm of $g$ restricted to a small set is small). This
  bound is \emph{uniform} in~$n$---it depends only on the $L^2$ bound
  on~$f_n$, not on the specific function.

\item \textbf{Vitali's convergence theorem:} on a finite measure space,
  if a sequence is uniformly integrable and converges a.e., then it
  converges in $L^1$. Applied here: $\int_K |(f_n - f)g| \to 0$.
\end{enumerate}

Combining: $|\int (f_n - f)g| < \varepsilon/3 + 2\varepsilon/3 =
\varepsilon$ for $n$ large enough.

\section{Where you'd get stuck}

\begin{itemize}
\item \textbf{Trying to use the Dominated Convergence Theorem directly.}
  You need a \emph{fixed} integrable function that dominates $|f_n \cdot
  g|$ for all~$n$. You don't have one. The $L^2$ bound gives
  $\|f_n\|_2 \le 1$, not $|f_n(x)| \le M$. This is the fundamental
  difference between $L^2$ control (average) and $L^\infty$ control
  (pointwise).
\item \textbf{The AM-GM trap.} $|f_n g| \le (f_n^2 + g^2)/2$ looks
  like a dominator, but $f_n^2$ depends on~$n$. A dominator must be
  fixed.
\item \textbf{Confusing weak and strong convergence.} $\int f_n g \to
  \int fg$ does NOT mean $\|f_n - f\|_2 \to 0$. The sequence
  $f_n(x) = \sin(nx)$ on $[0, 2\pi]$ converges to zero weakly
  (Riemann--Lebesgue) but $\|f_n\|_2 = \sqrt{\pi}$ for all~$n$.
\end{itemize}

% --------------------------------------------------------------------
\chapter{C3: Algebra --- Why $K[x,y]$ Can't Be a PID}\label{ch:lesson-c3}

\section{The setup}

A \textbf{PID} (Principal Ideal Domain) is a ring where every ideal is
generated by a single element. The integers $\mathbb{Z}$ are a PID: every
ideal looks like $(n) = \{nk : k \in \mathbb{Z}\}$. The polynomial ring
$K[x]$ (polynomials in one variable over a field~$K$) is also a PID:
you can always do polynomial long division, so the Euclidean algorithm
works, so every ideal is principal.

The question: what about $K[x, y]$, polynomials in \emph{two} variables?

\section{The key insight: one variable too many}

$K[x, y]$ is the same ring as $K[x][y]$---polynomials in~$y$ whose
coefficients are polynomials in~$x$. If $K[x][y]$ were a PID, what
would that say about~$K[x]$?

Here's the structural fact: \textbf{if $R$ is a PID, then $R[y]$ is a
PID if and only if $R$ is a field.} The ``only if'' direction is what
we need. If $R[y]$ is a PID and $R$ is a PID, then $R$ must be a field.

We know $K[x]$ is a PID. If $K[x][y]$ were also a PID, then $K[x]$
would be a field. But $K[x]$ is \emph{not} a field: the polynomial~$x$
has no multiplicative inverse (there's no polynomial $p(x)$ with
$x \cdot p(x) = 1$).

Contradiction. $K[x, y]$ is not a PID.

\section{What part (a) is really about}

The exam also asks you to prove that in an integral domain~$R$, if $a$
is nonzero and noninvertible, then the ideal $(a, x) \subset R[x]$ is
not principal. This is the engine behind part~(b): in $K[x][y]$,
take $R = K[x]$ and $a = x$. Then $(x, y)$ is a non-principal ideal
in $K[x, y]$.

The direct argument for~(a): suppose $(a, x) = (d)$ for some $d \in R[x]$.
Since $a \in (d)$ and $a$ is a constant (degree~$0$), $d$ must also be a
constant. Since $x \in (d)$ and $d$ is a constant, $d$ divides~$x$ in
$R[x]$, so $d$ is a unit. But then $(d) = R[x]$, while $(a, x) \ne R[x]$
because evaluating at $x = 0$ sends $(a, x)$ to $(a) \ne R$ (since $a$
is not a unit). Contradiction.

\section{Where you'd get stuck}

\begin{itemize}
\item \textbf{Not knowing the structure theorem.} If you don't know that
  ``$R[y]$ PID implies $R$ field'' (for $R$ a PID), you're stuck doing
  the generator-and-divisibility argument from scratch.
\item \textbf{Confusing ``$K[x]$ is a PID'' with ``$K[x]$ is a field.''}
  Being a PID means every ideal is principal. Being a field means every
  nonzero element has an inverse. $K[x]$ is the first but not the
  second. The exam is testing whether you understand this distinction.
\end{itemize}

% --------------------------------------------------------------------
\chapter{C4: Functional Analysis --- The Adjoint Trick}\label{ch:lesson-c4}

\section{The setup}

A \textbf{Hilbert space} is a complete inner product space: a vector space
with an inner product $\langle \cdot, \cdot \rangle$ and no ``missing
points'' (every Cauchy sequence converges). Think of $L^2(\mathbb{R})$
from C2, or $\mathbb{R}^n$ with the dot product, or $\ell^2$ (sequences
$(a_1, a_2, \ldots)$ with $\sum |a_n|^2 < \infty$).

A \textbf{linear map} $A : X \to X$ is a function respecting addition and
scalar multiplication. It is \textbf{continuous} (equivalently,
\textbf{bounded}) if there exists $C$ such that $\|Ax\| \le C\|x\|$ for
all~$x$. Linear maps on finite-dimensional spaces are always continuous.
On infinite-dimensional spaces, they need not be---which is what makes
this problem interesting.

The problem gives you $A$ (linear, not assumed continuous) and $B$ (not
assumed linear OR continuous!) satisfying $\langle Ax, y \rangle = \langle
x, By \rangle$ for all $x, y$. This is the \textbf{adjoint identity}: $B$
is behaving like the adjoint of~$A$, even though we haven't assumed it has
any nice properties.

\section{The tool you need: the Closed Graph Theorem}

\textbf{Theorem} (Banach, 1932). Let $X$ and $Y$ be Banach spaces
(complete normed spaces---Hilbert spaces are a special case) and let
$T : X \to Y$ be a linear map. If the graph of~$T$,
\[
  \mathrm{graph}(T) = \{(x, Tx) : x \in X\} \subseteq X \times Y,
\]
is a \emph{closed set}, then $T$ is continuous.

``Closed'' means: if $(x_n, Tx_n) \to (x, z)$ in $X \times Y$, then
$z = Tx$. In other words, if the inputs converge and the outputs converge,
then the outputs converge to the right thing.

This is a remarkable theorem because it lets you prove continuity
\emph{without ever bounding $\|Tx\|$}. You just check that limits
behave consistently.

\section{Part (a): B is linear --- nondegeneracy does the work}

The inner product is \textbf{nondegenerate}: if $\langle x, a \rangle =
\langle x, b \rangle$ for all~$x$, then $a = b$. (Proof: take $x =
a - b$, get $\|a - b\|^2 = 0$.)

To show $B(y_1 + y_2) = By_1 + By_2$: for all~$x$,
\begin{align*}
  \langle x, B(y_1 + y_2) \rangle
  &= \langle Ax, y_1 + y_2 \rangle
  = \langle Ax, y_1 \rangle + \langle Ax, y_2 \rangle \\
  &= \langle x, By_1 \rangle + \langle x, By_2 \rangle
  = \langle x, By_1 + By_2 \rangle.
\end{align*}
By nondegeneracy, $B(y_1 + y_2) = By_1 + By_2$. Same argument for
$B(cy) = cBy$, using linearity of the inner product in the second
argument.

\section{Parts (b) and (c): the Closed Graph Theorem}

Now that $B$ is linear, we can apply the Closed Graph Theorem.

\textbf{Claim:} $A$ has closed graph.

\textbf{Proof of claim:} Suppose $x_n \to x$ and $Ax_n \to z$. We
need $z = Ax$. For any $w \in X$:
\begin{align*}
  \langle w, z \rangle
  &= \lim_n \langle w, Ax_n \rangle
  &\text{(inner product is continuous in the second variable)} \\
  &= \lim_n \overline{\langle Ax_n, w \rangle}
  &\text{(conjugate symmetry: $\langle a, b \rangle = \overline{\langle b, a \rangle}$)} \\
  &= \lim_n \overline{\langle x_n, Bw \rangle}
  &\text{(adjoint identity with $y = w$)} \\
  &= \overline{\langle x, Bw \rangle}
  &\text{(inner product is continuous in the first variable)} \\
  &= \langle Bw, x \rangle
  &\text{(conjugate symmetry again)} \\
  &= \overline{\langle x, Bw \rangle} = \langle Bw, x \rangle \\
  &\quad\;\;\; \text{wait---let me redo this without the double flip} \\[4pt]
  &= \lim_n \langle Bw, x_n \rangle
  &\text{(conjugate symmetry applied to $\langle x_n, Bw \rangle$)} \\
  &= \langle Bw, x \rangle = \langle w, Ax \rangle.
  &\text{(adjoint identity: $\langle Bw, x \rangle = \overline{\langle x, Bw \rangle} = \overline{\overline{\langle Bw, x \rangle}} \ldots$)}
\end{align*}

OK---the conjugate symmetry bookkeeping is genuinely fiddly, even
informally. Let me state the clean version:

\begin{quote}
For real Hilbert spaces (no conjugates): $\langle w, Ax_n \rangle =
\langle Bw, x_n \rangle$ by the adjoint identity with the variables
swapped. Taking limits: $\langle w, z \rangle = \langle Bw, x \rangle
= \langle w, Ax \rangle$. By nondegeneracy, $z = Ax$.
\end{quote}

For complex Hilbert spaces, you have to route through two applications of
$\langle a, b \rangle = \overline{\langle b, a \rangle}$ to get from
$\langle w, Ax_n \rangle$ to $\langle Bw, x_n \rangle$. This is exactly
what tripped up the Lean formalisation---three failed tactic attempts
before the right rewrite chain was found.

By the Closed Graph Theorem, $A$ is continuous. For~$B$: the adjoint
identity is symmetric. $\langle Bx, y \rangle = \overline{\langle y,
Bx \rangle} = \overline{\langle Ay, x \rangle} = \langle x, Ay \rangle$.
So $B$ has adjoint~$A$, and the same closed-graph argument gives~$B$
continuous.

\section{Where you'd get stuck}

\begin{itemize}
\item \textbf{Not recognising the Closed Graph Theorem setup.} The
  problem says ``do not assume continuous.'' This is the exam waving a
  flag: \emph{you will need a theorem that gives you continuity for free
  from some algebraic condition.} In a Hilbert/Banach space, that
  theorem is the Closed Graph Theorem (or equivalently, the Open Mapping
  Theorem).
\item \textbf{Trying to bound $\|Ax\|$ directly.} You can't---the
  whole point is that you don't have a bound. The Closed Graph Theorem
  lets you avoid this entirely.
\item \textbf{The conjugate symmetry dance.} In a complex Hilbert space,
  $\langle a, b \rangle = \overline{\langle b, a \rangle}$, which means
  the inner product is linear in one argument and \emph{conjugate}-linear
  in the other. Swapping the arguments of the adjoint identity requires
  tracking these conjugates. This is where both humans and proof
  assistants stumble. The Lean formalisation needed explicit
  \texttt{inner\_conj\_symm} rewrites that the informal proof waves
  away.
\end{itemize}


% ====================================================================
\part{The Proofs}

% --------------------------------------------------------------------
\chapter{C1: Topology}\label{ch:c1}

\begin{problem}[UT Austin, t01A01, 2001]
Let $f : \RP^2 \to \RP^2$ be a continuous map that is not surjective.
\begin{enumerate}[label=(\alph*)]
\item Show that $f_* : \pi_1(\RP^2) \to \pi_1(\RP^2)$ has image $\{1\}$.
\item Show that $f$ is homotopic to a constant map.
\end{enumerate}
\emph{Hint: What does $\RP^2 \setminus \{\mathrm{point}\}$ look like?}
\end{problem}

\begin{proofbox}
\textbf{(a)} Since $f$ is not surjective, it misses some point $q \in \RP^2$.
So $f$ factors through $\RP^2 \setminus \{q\}$, which is homeomorphic to an
open M\"obius band and therefore has the homotopy type of~$S^1$. Thus
$\pi_1(\RP^2 \setminus \{q\}) \cong \mathbb{Z}$.

The induced map on fundamental groups factors as
\[
  \pi_1(\RP^2) \xrightarrow{g_*} \pi_1(\RP^2 \setminus \{q\})
  \xrightarrow{i_*} \pi_1(\RP^2),
\]
i.e.\ $\mathbb{Z}/2\mathbb{Z} \xrightarrow{g_*} \mathbb{Z}
\xrightarrow{i_*} \mathbb{Z}/2\mathbb{Z}$.

The only group homomorphism $\mathbb{Z}/2\mathbb{Z} \to \mathbb{Z}$
is trivial, because $\mathbb{Z}$ is torsion-free: if $\varphi(1) = k$,
then $0 = \varphi(1+1) = 2k$, so $k = 0$. Therefore $g_* = 0$, hence
$f_* = i_* \circ g_* = 0$.

\medskip
\textbf{(b)} Since $f_*$ is trivial, $f$ lifts to a map
$\tilde{f} : \RP^2 \to S^2$ (the universal cover). The lift
$\tilde{f}$ also misses a point (the preimage of~$q$ under the
covering map includes at least one point not in the image of
$\tilde{f}$). So $\tilde{f}$ factors through $S^2 \setminus \{p\}
\cong \mathbb{R}^2$, which is contractible. Hence $\tilde{f}$ is
nullhomotopic, and therefore $f = p \circ \tilde{f}$ is nullhomotopic.
\end{proofbox}

\begin{leanbox}
\begin{Verbatim}[fontsize=\small]
/-- The only group homomorphism Z/2Z -> Z is trivial. -/
theorem hom_zmod2_to_int_trivial
    (phi : ZMod 2 ->+ Z) : phi = 0 := by
  ext x
  fin_cases x
  . exact map_zero phi
  . have h2 : (1 : ZMod 2) + 1 = 0 := by decide
    have h := phi.map_add 1 1
    rw [h2, map_zero] at h
    have key : phi 1 + phi 1 = 0 := h.symm
    have : (2 : Z) * phi 1 = 0 := by ring_nf; linarith
    have := mul_eq_zero.mp this
    rcases this with h2 | done
    . norm_num at h2
    . exact done

/-- Any composite Z/2Z ->+ Z ->+ Z/2Z is trivial. -/
theorem composite_through_Z_trivial
    (phi : ZMod 2 ->+ Z) (psi : Z ->+ ZMod 2) :
    psi.comp phi = 0 := by
  rw [hom_zmod2_to_int_trivial phi]; simp
\end{Verbatim}
\end{leanbox}

\begin{notebox}
The algebraic core---the only homomorphism
$\mathbb{Z}/2\mathbb{Z} \to \mathbb{Z}$ is trivial---is fully verified.
The topological setup ($\RP^2$, covering spaces, homotopy) is not in
Mathlib and is stated as \texttt{True := trivial} to mark the boundary.
The Lean proof of the algebraic lemma is close to the informal one:
\texttt{fin\_cases} dispatches the two elements of $\mathbb{Z}/2\mathbb{Z}$,
\texttt{map\_add} applies the homomorphism property, and the conclusion
$2k = 0 \Rightarrow k = 0$ is handled by \texttt{ring\_nf + linarith + mul\_eq\_zero}.
\end{notebox}

% --------------------------------------------------------------------
\chapter{C2: Analysis}\label{ch:c2}

\begin{problem}[UT Austin, a96A03, 1996]
Let $B = \{f \in L^2(\mathbb{R}) : \|f\|_2 \le 1\}$ and let $\{f_n\}$
be a sequence from~$B$ with $f_n \to f$ almost everywhere. Show that:
\begin{enumerate}[label=(\roman*)]
\item $f \in B$\,;
\item $\displaystyle\lim_{n \to \infty} \int_{-\infty}^{\infty}
  f_n(x)\,g(x)\,dx = \int_{-\infty}^{\infty} f(x)\,g(x)\,dx$
  for all $g \in L^2(\mathbb{R})$.
\end{enumerate}
\emph{Hint: for the second part, start by considering the case in which
$g$ is continuous and supported in a bounded interval.}
\end{problem}

\begin{proofbox}
\textbf{(i)} By Fatou's lemma applied to $|f_n|^2$:
\[
  \|f\|_2^2 = \int |f|^2 \le \liminf_n \int |f_n|^2
  = \liminf_n \|f_n\|_2^2 \le 1.
\]

\textbf{(ii)} We show $\int (f_n - f)\,g \to 0$. Fix $\varepsilon > 0$.

\emph{Tail control.} Choose a measurable set $s$ with $\mu(s) < \infty$
such that $\|s^c\text{-indicator of } g\|_2 < \varepsilon/3$. This is
possible because $g \in L^2$.

\emph{Core convergence.} On $s$ (finite measure), the sequence
$\{(f_n - f) \cdot g\}$ is $L^1$-uniformly integrable. The key
estimate: for any measurable $A \subseteq s$ with $\mu(A) \le \delta$,
H\"older gives
\[
  \left\| \mathbf{1}_A \cdot (f_n - f) \cdot g \right\|_1
  \le \|f_n - f\|_2 \cdot \left\| \mathbf{1}_A \cdot g \right\|_2
  \le 2 \cdot \left\| \mathbf{1}_A \cdot g \right\|_2,
\]
and the absolute continuity of the $L^2$ norm of~$g$ makes
$\|\mathbf{1}_A \cdot g\|_2$ small when $\mu(A)$ is small (uniformly
in~$n$). By the Vitali convergence theorem ($L^1$ on a finite measure),
$\int_s (f_n - f)\,g \to 0$.

\emph{Assembly.} For $n \ge N$:
\begin{align*}
  \left| \int (f_n - f)\,g \right|
  &\le \left| \int_s (f_n - f)\,g \right|
     + \left| \int_{s^c} (f_n - f)\,g \right| \\
  &< \frac{\varepsilon}{3}
     + \|f_n - f\|_2 \cdot \|s^c\text{-indicator of } g\|_2 \\
  &\le \frac{\varepsilon}{3} + 2 \cdot \frac{\varepsilon}{3}
  = \varepsilon.
\end{align*}
\end{proofbox}

\begin{leanbox}
\begin{Verbatim}[fontsize=\scriptsize]
-- Part (i): one-liner using Mathlib's Fatou-for-eLpNorm.
theorem ae_limit_in_L2_ball ... :=
  MeasureTheory.Lp.eLpNorm_le_of_ae_tendsto
    (Eventually.of_forall hF_bound) hF_meas hae

-- Part (ii): 320+ lines assembling:
-- 1. abs_integral_mul_le_eLpNorm (Cauchy-Schwarz for integrals)
-- 2. unifIntegrable_one_of_L2_bound (L1-UI from L2 bound via Hoelder)
-- 3. unifIntegrable_product (UI of (Fn-f)*g from abs. continuity of g)
-- 4. tendsto_setIntegral_mul (L1-Vitali on finite-measure restriction)
-- 5. weak_convergence_L2 (epsilon/3 tail + core assembly)
--
-- Key Mathlib APIs used:
--   eLpNorm_smul_le_mul_eLpNorm, eLpNorm_le_eLpNorm_mul_rpow_measure_univ,
--   tendsto_Lp_finite_of_tendsto_ae, MemLp.eLpNorm_indicator_le,
--   MemLp.exists_eLpNorm_indicator_compl_lt, integral_add_compl,
--   tendsto_integral_of_L1'
\end{Verbatim}
\end{leanbox}

\begin{notebox}
C2 was the hardest canary. The informal proof follows the textbook
hint---start with compactly supported~$g$, then extend by density---but
the formalisation takes a slightly different route: uniform integrability
of the \emph{product} $(f_n - f) \cdot g$ rather than of~$f_n$ alone,
because this lets the absolute continuity of $g$'s $L^2$ norm do the work
that a pointwise dominator cannot (we have no $L^\infty$ bound on~$f_n$).

The 356-line proof is almost entirely API composition: every mathematical
step corresponds to a named Mathlib lemma. The difficulty was not missing
mathematics but type-coercion plumbing---converting between
\texttt{eLpNorm}, \texttt{ENNReal}, \texttt{ofReal}, indicator functions,
and measure restrictions. The informal proof is five lines; the formal
one is seventy times longer. This ratio is the current cost of precision.
\end{notebox}

% --------------------------------------------------------------------
\chapter{C3: Algebra}\label{ch:c3}

\begin{problem}[UT Austin, b03J02, 2003]
Let $R$ be an integral domain and let $a$ be a nonzero noninvertible
element of~$R$.
\begin{enumerate}[label=(\alph*)]
\item Prove that the ideal $(a, x) \subset R[x]$ is not principal.
\item Use~(a) to show that if $K$ is a field then $K[x, y]$ is not a PID.
\end{enumerate}
\end{problem}

\begin{proofbox}
\textbf{(b)} We have $K[x, y] \cong K[x][y]$. If $K[x][y]$ were a PID,
then $K[x]$ would be a field (since for a PID~$R$, $R[y]$ is a PID if
and only if~$R$ is a field). But $K[x]$ is not a field: $x$ is
irreducible but not a unit. Contradiction.

\medskip
\textbf{(a)} is the lemma used above: in $R[x]$, the ideal $(a, x)$ is
not principal because any generator~$d$ would divide both~$a$ (a
constant) and~$x$, forcing $d$ to be a unit, but then $(d) = R[x]$
while $(a, x) \ne R[x]$ (since $a$ is not a unit, the evaluation
at~$x = 0$ sends $(a, x)$ into~$(a) \ne R$).
\end{proofbox}

\begin{leanbox}
\begin{Verbatim}[fontsize=\small]
/-- K[X][Y] is not a PID. -/
theorem kxy_not_pid (K : Type*) [Field K] [DecidableEq K] :
    ¬ IsPrincipalIdealRing (Polynomial (Polynomial K)) := by
  intro h
  have hfield : IsField K[X] :=
    IsField.of_isPrincipalIdealRing_polynomial (R := K[X])
  exact polynomial_not_isField hfield
\end{Verbatim}
\end{leanbox}

\begin{notebox}
The shortest canary: 19~lines including comments. Mathlib already knows
that $K[X]$ is a Euclidean domain (hence PID) but not a field
(\texttt{polynomial\_not\_isField}), and that if~$R[Y]$ is a PID
then~$R$ is a field
(\texttt{IsField.of\_isPrincipalIdealRing\_polynomial}). The
formalisation is a two-step contradiction.

Part~(a)---the non-principality of $(a, x)$---is the lemma that
Mathlib already encapsulates. The exam splits it out as a separate
subproblem; the formalisation collapses it into the library.
\end{notebox}

% --------------------------------------------------------------------
\chapter{C4: Functional Analysis}\label{ch:c4}

\begin{problem}[UT Austin, m96A04, 1996]
Let $X$ be a Hilbert space and let $A : X \to X$ be a linear map (not
assumed continuous). Let $B : X \to X$ satisfy
$\langle Ax, y \rangle = \langle x, By \rangle$ for all $x, y \in X$.
(Do not assume that $B$ is linear or continuous.) Prove:
\begin{enumerate}[label=(\alph*)]
\item $B$ is linear.
\item $B$ is continuous.
\item $A$ is continuous.
\end{enumerate}
\end{problem}

\begin{proofbox}
\textbf{(a)} Fix $y_1, y_2 \in X$ and $c \in \mathbb{C}$. For all $x$:
\[
  \langle x, B(y_1 + y_2) \rangle = \langle Ax, y_1 + y_2 \rangle
  = \langle Ax, y_1 \rangle + \langle Ax, y_2 \rangle
  = \langle x, By_1 \rangle + \langle x, By_2 \rangle
  = \langle x, By_1 + By_2 \rangle.
\]
By nondegeneracy of the inner product, $B(y_1 + y_2) = By_1 + By_2$.
Similarly $B(cy) = cBy$.

\medskip
\textbf{(b) and (c)} By the Closed Graph Theorem. We show $A$ has
closed graph; the argument for~$B$ is symmetric.

Suppose $x_n \to x$ and $Ax_n \to z$. For all $w$:
\[
  \langle w, z \rangle = \lim_n \langle w, Ax_n \rangle
  = \lim_n \langle Bw, x_n \rangle = \langle Bw, x \rangle
  = \langle w, Ax \rangle.
\]
(The second equality uses the adjoint identity rewritten via conjugate
symmetry: $\langle w, Ax_n \rangle = \overline{\langle Ax_n, w \rangle}
= \overline{\langle x_n, Bw \rangle} = \langle Bw, x_n \rangle$.)

By nondegeneracy, $z = Ax$. So the graph of~$A$ is sequentially closed,
hence closed (in a Banach space). By the Closed Graph Theorem, $A$ is
continuous.

For~$B$: the adjoint identity gives $\langle Bx, y \rangle =
\overline{\langle y, Bx \rangle} = \overline{\langle Ay, x \rangle}
= \langle x, Ay \rangle$, so $B$ has adjoint~$A$ and the same argument
applies.
\end{proofbox}

\begin{leanbox}
\begin{Verbatim}[fontsize=\small]
-- (a): B is additive and respects scalars.
theorem adjoint_map_additive ... := by
  intro y1 y2
  refine ext_inner_left C fun x => ?_
  simp only [<- hadj, inner_add_right]

-- (b,c): helper proves the closed-graph condition.
theorem adjoint_seq_closed ...
    (hu : Tendsto u atTop (nhds x))
    (hAu : Tendsto (A . u) atTop (nhds z)) :
    z = A x := by
  refine ext_inner_left C fun w => ?_
  -- w,z = lim w,A(un) = lim Bw,un = Bw,x = w,Ax
  have h1 := (continuous_const.inner continuous_id').tendsto z
              |>.comp hAu  -- lim <w, A(un)> = <w, z>
  have h2 : ...            -- lim <w, A(un)> = <w, Ax>
    -- (uses inner_conj_symm to flip the adjoint identity)
  exact tendsto_nhds_unique h1 h2

-- (c): A is continuous by Closed Graph Theorem.
theorem adjoint_map_A_continuous ... : Continuous A :=
  A.continuous_of_seq_closed_graph
    (adjoint_seq_closed A B hadj)

-- (b): B is continuous (symmetric argument).
theorem adjoint_map_B_continuous ... : Continuous B := by
  apply B.continuous_of_seq_closed_graph
  exact adjoint_seq_closed B A hadj' u x z hu hBu
\end{Verbatim}
\end{leanbox}

\begin{notebox}
The informal proof has a natural three-act structure: linearity from
nondegeneracy, then the closed-graph argument, then the symmetry for~$B$.
The formal proof follows the same structure.

The nondegeneracy step uses \texttt{ext\_inner\_left}---Mathlib's name
for the fact that $\langle x, a \rangle = \langle x, b \rangle$ for all
$x$ implies $a = b$. The closed-graph step uses
\texttt{LinearMap.continuous\_of\_seq\_closed\_graph}, which is Mathlib's
sequentially-closed-graph variant of Banach's theorem.

The one subtlety: flipping the adjoint identity for~$B$ requires
navigating conjugate symmetry ($\langle a, b \rangle = \overline{\langle
b, a \rangle}$). The informal proof does this in a parenthetical; the
formal proof requires explicit \texttt{inner\_conj\_symm} rewrites with
careful attention to which variable is which. This was the main source
of failed tactic applications during development.
\end{notebox}

% ====================================================================
\part{Retrospective: The Discipline}

\chapter{M-apm-solutions and the Proof Peripheral}

\section{What this report is really about}

This report is not ``Claude solving maths problems.'' It is a canary
run for \textbf{M-apm-solutions}, a FUTON stack mission to produce
expository solutions to 489~UT Austin preliminary exam problems using a
disciplined, multi-system pipeline. The canaries test whether the
pipeline's components---the proof peripheral's phase discipline, the
question-asking patterns, Lean formalisation, and pedagogical
rotation---work when composed end-to-end.

The distinction matters. A raw language model can solve prelim problems
and explain them. What it cannot do, without external discipline, is
produce outputs with a consistent pedagogical structure, document its
dead ends, or separate what it proved from what it assumed. The
discipline provides the structure. The model provides the content.

\section{The proof peripheral's phase discipline}

The FUTON proof peripheral (\texttt{futon3c/peripheral/proof.clj})
structures every proof attempt as a sequence of phases:

\begin{center}
\begin{tabular}{lll}
\toprule
\textbf{Phase} & \textbf{Proof use} & \textbf{Pedagogical rotation} \\
\midrule
Observe  & Survey problem, identify blockers & ``What's really being asked'' \\
Propose  & Suggest approach with success criteria & ``The key insight'' \\
Execute  & Run the approach & ``The clean proof'' \\
Validate & Check artifacts against criteria & ``Does the argument close?'' \\
Classify & Status transition & Grade the proof \\
Integrate & Merge into ledger & --- \\
Commit   & Finalise, record evidence & --- \\
Gate-review & futon3b validation & --- \\
\bottomrule
\end{tabular}
\end{center}

The first four phases map directly to the lesson structure in Part~I.
Each chapter in ``The Lessons'' is organised as:
\textbf{setup} (Observe), \textbf{key insight} (Propose),
\textbf{proof} (Execute), \textbf{where you'd get stuck} (Validate).
This structure was not improvised during writing---it was specified in
the M-apm-solutions mission document before any mathematics was attempted.

\section{The question-asking patterns}

The FUTON stack maintains a set of question-asking patterns (QP-1 through
QP-8) calibrated on research mathematics. M-apm-solutions defines a
\textbf{pedagogical rotation} of five of these:

\begin{center}
\begin{tabular}{lp{4cm}p{5cm}l}
\toprule
\textbf{Type} & \textbf{QP source} & \textbf{Teaching form} & \textbf{Phase} \\
\midrule
\texttt{why-hard} & QP-4 (Failure Characterisation) &
  ``Why is this problem hard? What trap does it set?'' & Observe \\
\texttt{what-crux} & QP-7 (Kernel Identification) &
  ``What is the key insight? What do you need to see?'' & Propose \\
\texttt{why-works} & QP-5 (Theorem Applicability) &
  ``Why does this approach work? What hypotheses are being used?'' & Validate \\
\texttt{what-connects} & QP-1 (Landscape Scout) &
  ``What does this connect to? Where else does this pattern appear?'' & Integrate \\
\texttt{confidence} & QP-8 (Confidence Inversion) &
  ``What is surprising here? Where is intuition wrong?'' & Integrate \\
\bottomrule
\end{tabular}
\end{center}

Every ``Where you'd get stuck'' section in Part~I was generated by
applying \texttt{why-hard} (QP-4). Every ``key insight'' section was
generated by \texttt{what-crux} (QP-7). The ``dead end that teaches
something'' notes are \texttt{confidence} (QP-8). The patterns produce
the structure; the mathematical content fills it.

\section{The failed-route ledger}

The mission spec says: ``The failed-route ledger is pedagogically
gold---dead ends are where intuition lives.'' In the canary run:

\begin{description}
\item[C1:] Trying to show $\RP^2 \setminus \{q\}$ is contractible
  (it retracts to~$S^1$). This is a \texttt{confidence} (QP-8)
  moment: the intuition that ``removing a point simplifies'' is right,
  but ``simplifies'' means $\pi_1 = \mathbb{Z}$, not $\pi_1 = 0$.
\item[C2:] The AM-GM dominator trap ($|f_n g| \le (f_n^2 + g^2)/2$
  depends on~$n$). This is a \texttt{why-hard} (QP-4) moment: the
  problem is hard precisely because there's no pointwise dominator.
\item[C4:] Three failed \texttt{rw [inner\_conj\_symm]} attempts in
  Lean. This is a \texttt{why-works} (QP-5) moment: the conjugate
  symmetry that the informal proof hand-waves requires explicit tracking
  of which variable is conjugated.
\item[C3:] No significant dead ends---the proof was too short.
\end{description}

In a raw Claude session without the ledger discipline, these dead ends
would be silently discarded. The output would be a clean proof with no
record of what was tried and failed. The pedagogical value---the
intuition-building---would be lost.

\section{The canary profiles}

Each canary was chosen to stress-test a different aspect of the pipeline:

\begin{center}
\begin{tabular}{lllp{5.5cm}}
\toprule
\textbf{Canary} & \textbf{Subject} & \textbf{Lean lines} & \textbf{What it tests} \\
\midrule
C1 & Topology   & 48  & Mathlib boundary (missing coverage) \\
C2 & Analysis   & 356 & API composition at scale \\
C3 & Algebra    & 19  & Library-level proof (near-free) \\
C4 & Functional & 94  & Named-theorem application + conjugate bookkeeping \\
\bottomrule
\end{tabular}
\end{center}

The 489-problem pipeline will encounter all four profiles repeatedly.
C3-type problems (the library already knows) should be automated.
C1-type problems (Mathlib doesn't cover the topology) need the
boundary marked cleanly. C4-type problems (a named theorem plus
careful plumbing) are the sweet spot. C2-type problems (massive API
composition) are where the pipeline will spend most of its time.

\section{What the discipline contributed}

The natural question is: how much of this would a raw Claude session
produce without the FUTON discipline?

The proofs: probably equivalent. The mathematical content comes from the
model's training, not from the proof peripheral. A clean prompt---``solve
these four problems and formalise them in Lean''---would likely produce
similar Lean code.

The pedagogy: probably thinner. The five-question structure
(\texttt{why-hard}, \texttt{what-crux}, \texttt{why-works},
\texttt{what-connects}, \texttt{confidence}) is not the model's default
register. Without the phase discipline, a raw session would produce
``here's the proof, now here's what it means''---linear exposition rather
than structured tutoring. The ``where you'd get stuck'' sections and the
dead-end analysis would likely be absent or shallow.

The failed-route documentation: almost certainly absent. Language models
do not spontaneously record what they tried and discarded. The ledger
discipline forces this.

A \textbf{control experiment}---giving a fresh Claude session the same
problems with a generic prompt---would make this comparison empirical
rather than speculative. It is planned as a companion piece.

\section{Relation to the First Proof monograph}

The First Proof sprint~\cite{firstproof} attempted ten open research
problems in 55~hours. The canaries attempted four textbook problems in a
single session. The difference in ambition reflects a difference in
purpose: the First Proof sprint was exploratory (``what can we
attempt?''); the canaries are diagnostic (``what does the pipeline
need?'').

The methodological continuity is in the discipline: informal proof first,
then formalisation, with the formal step exposing gaps and subtleties
that the informal argument conceals. The Lakatosian
rhythm---conjecture, critique, repair---operates here at the level of
tactic application rather than theorem statement. The proof peripheral's
phases are the same; the question-asking patterns are the same; the
failed-route ledger is the same. What changes is the scale (4~problems
instead of~10) and the formality (Lean~4 instead of prose-level
verification).

The deeper continuity is architectural. The First Proof sprint produced
the proof peripheral and the question patterns. The canaries are the
first test of whether those tools transfer from research mathematics to
expository mathematics. They do. The 489-problem pipeline is the next
step: the same discipline, applied at scale, with Arxana retaining the
full graph and Tickle dispatching the automation.

\begin{thebibliography}{9}
\bibitem{firstproof}
  J.~Corneli. \emph{A First Proof Sprint}. Hyperreal Enterprises, 2026.
  \url{https://github.com/tothedarktowercame/futon6}

\bibitem{lakatos1976}
  I.~Lakatos. \emph{Proofs and Refutations}. Cambridge University Press, 1976.

\bibitem{pease2017}
  A.~Pease, J.~Corneli, et al.
  Lakatos-style collaborative mathematics through dialectical, structured
  and abstract argumentation.
  \emph{Artificial Intelligence}, 246:181--219, 2017.
\end{thebibliography}

\end{document}
